% oss-plantilla: oss-2026.11
% Preámbulo corporativo fijo. No editar en sd-NN.tex: exportar_latex.py verificar lo compara byte a byte.
% Para cambiar el formato, editar esta plantilla u oss.sty y subir la versión de la primera línea.
\PassOptionsToPackage{unicode}{hyperref}
\PassOptionsToPackage{hyphens}{url}
\documentclass[11pt]{article}
\usepackage{amsmath,amssymb}
\usepackage{iftex}
\ifPDFTeX
  % Prism y otros editores en la nube: pdfLaTeX con Helvetica como familia por defecto.
  \usepackage[T1]{fontenc}
  \usepackage[utf8]{inputenc}
  \usepackage{textcomp}
  \usepackage[scaled=0.92]{helvet}
  \renewcommand{\familydefault}{\sfdefault}
  \DeclareUnicodeCharacter{2264}{\ensuremath{\leq}}
  \DeclareUnicodeCharacter{2265}{\ensuremath{\geq}}
  \DeclareUnicodeCharacter{2260}{\ensuremath{\neq}}
  \DeclareUnicodeCharacter{2248}{\ensuremath{\approx}}
  \DeclareUnicodeCharacter{2190}{\ensuremath{\leftarrow}}
  \DeclareUnicodeCharacter{2192}{\ensuremath{\rightarrow}}
  \DeclareUnicodeCharacter{25CF}{\ensuremath{\bullet}}
  \DeclareUnicodeCharacter{2713}{\ensuremath{\checkmark}}
\else
  % XeLaTeX/LuaLaTeX (compilación oficial local): DejaVu Sans.
  \usepackage{fontspec}
  \defaultfontfeatures{Ligatures=TeX,Scale=MatchLowercase}
  \setmainfont{DejaVu Sans}
  \setsansfont{DejaVu Sans}
  \setmonofont{DejaVu Sans Mono}
\fi
\usepackage[letterpaper,margin=20mm]{geometry}
\usepackage{microtype}
\usepackage{parskip}
\usepackage{xcolor}
\usepackage{longtable,booktabs,array}
\usepackage{calc}
\usepackage{etoolbox}
\makeatletter
\patchcmd\longtable{\par}{\if@noskipsec\mbox{}\fi\par}{}{}
\makeatother
\usepackage{footnotehyper}
\makesavenoteenv{longtable}
\usepackage{graphicx}
\makeatletter
\newsavebox\pandoc@box
\newcommand*\pandocbounded[1]{%
  \sbox\pandoc@box{#1}%
  \Gscale@div\@tempa{\textheight}{\dimexpr\ht\pandoc@box+\dp\pandoc@box\relax}%
  \Gscale@div\@tempb{\linewidth}{\wd\pandoc@box}%
  \ifdim\@tempb\p@<\@tempa\p@\let\@tempa\@tempb\fi
  \ifdim\@tempa\p@<\p@\scalebox{\@tempa}{\usebox\pandoc@box}%
  \else\usebox{\pandoc@box}%
  \fi}
\def\maxwidth{\ifdim\Gin@nat@width>\linewidth\linewidth\else\Gin@nat@width\fi}
\def\maxheight{\ifdim\Gin@nat@height>\textheight\textheight\else\Gin@nat@height\fi}
\makeatother
\setkeys{Gin}{width=\maxwidth,height=\maxheight,keepaspectratio}
\makeatletter
\def\fps@figure{htbp}
\makeatother
\usepackage[normalem]{ulem}
\providecommand{\st}[1]{\sout{#1}}
\usepackage{fancyvrb}
\setlength{\emergencystretch}{3em}
\providecommand{\tightlist}{%
  \setlength{\itemsep}{0pt}\setlength{\parskip}{0pt}}
\providecommand{\passthrough}[1]{#1}
\setcounter{secnumdepth}{-\maxdimen}
\ifPDFTeX
  \usepackage[spanish]{babel}
\else
  \usepackage[bidi=default]{babel}
  \babelprovide[main,import]{spanish}
\fi
\let\LanguageShortHands\languageshorthands
\def\languageshorthands#1{}
\usepackage{bookmark}
\usepackage{xurl}
\urlstyle{same}
\hypersetup{pdflang={es-CL},pdfcreator={Only Simple Solutions - exportar_latex.py}}
\usepackage{oss}
\begin{document}

\ossCover{Especificaciones técnicas ofertadas}{Subdocumento 4 · T-11}
\tableofcontents
\clearpage
\section*{Lista de tablas}
\addcontentsline{toc}{section}{Lista de tablas}

\noindent\hangindent=2em\hyperref[tab:t11-1]{Tabla T-11.1: Correspondencia entre el formulario y la propuesta}\dotfill\pageref{tab:t11-1}\par\vspace{0.4\baselineskip}
\noindent\hangindent=2em\hyperref[tab:t11-2]{Tabla T-11.2: Servicios de Azure contratados}\dotfill\pageref{tab:t11-2}\par\vspace{0.4\baselineskip}
\noindent\hangindent=2em\hyperref[tab:t11-3]{Tabla T-11.3: Plataforma y software de la solución}\dotfill\pageref{tab:t11-3}\par\vspace{0.4\baselineskip}
\noindent\hangindent=2em\hyperref[tab:t11-4]{Tabla T-11.4: Herramientas de seguridad, entrega y operación}\dotfill\pageref{tab:t11-4}\par\vspace{0.4\baselineskip}
\noindent\hangindent=2em\hyperref[tab:t11-5]{Tabla T-11.5: Cómputo de tienda provisto por Only Simple Solutions}\dotfill\pageref{tab:t11-5}\par\vspace{0.4\baselineskip}
\noindent\hangindent=2em\hyperref[tab:t11-6]{Tabla T-11.6: Especificación de hardware y obras que adquiere el Cliente}\dotfill\pageref{tab:t11-6}\par\vspace{0.4\baselineskip}
\noindent\hangindent=2em\hyperref[tab:t11-7]{Tabla T-11.7: Enlaces por sitio}\dotfill\pageref{tab:t11-7}\par\vspace{0.4\baselineskip}
\noindent\hangindent=2em\hyperref[tab:t11-8]{Tabla T-11.8: Especificaciones de los centros de datos}\dotfill\pageref{tab:t11-8}\par\vspace{0.4\baselineskip}
\noindent\hangindent=2em\hyperref[tab:t11-dc01]{Tabla T-11.9: Equipos y obras de DC-01 y reparto de suministro}\dotfill\pageref{tab:t11-dc01}\par\vspace{0.4\baselineskip}
\noindent\hangindent=2em\hyperref[tab:t11-9]{Tabla T-11.10: Plataformas del Cliente y su tratamiento}\dotfill\pageref{tab:t11-9}\par\vspace{0.4\baselineskip}
\noindent\hangindent=2em\hyperref[tab:t11-10]{Tabla T-11.11: Componentes de las innovaciones}\dotfill\pageref{tab:t11-10}\par\vspace{0.4\baselineskip}
\noindent\hangindent=2em\hyperref[tab:t11-11]{Tabla T-11.12: Licenciamiento por producto}\dotfill\pageref{tab:t11-11}\par\vspace{0.4\baselineskip}
\noindent\hangindent=2em\hyperref[tab:t11-12]{Tabla T-11.13: Versión, fin de soporte y plan de actualización}\dotfill\pageref{tab:t11-12}\par\vspace{0.4\baselineskip}
\noindent\hangindent=2em\hyperref[tab:t11-13]{Tabla T-11.14: Portabilidad de los componentes}\dotfill\pageref{tab:t11-13}\par\vspace{0.4\baselineskip}
\noindent\hangindent=2em\hyperref[tab:t11-14]{Tabla T-11.15: Comprobaciones que fijan cantidades, niveles y modelos}\dotfill\pageref{tab:t11-14}\par\vspace{0.4\baselineskip}
\clearpage


\renewcommand{\tablename}{Tabla}
\renewcommand{\thetable}{T-11.\arabic{table}}

\section{Formulario T-11 · Especificaciones técnicas ofertadas}\label{formulario-t-11}

Conforme al Formulario T-11 de las Bases Administrativas, que se incluye en el Formulario T-7, este formulario detalla los componentes de infraestructura, plataforma, licenciamiento y hardware especificado de la solución ofertada a Multitiendas Ancoa S.A., con su ubicación y su justificación. Las secciones 4.1.1, 4.2, 4.2.1 y 4.3 del Subdocumento 4 contienen el resumen y el análisis. Este formulario reúne el inventario técnico, incluidos los componentes que el Subdocumento 4 menciona en el texto y en las figuras junto con el mapeo de la Tabla 4.20.

El reparto del suministro sigue la Tabla 4.22 del SD-04. Only Simple Solutions provee el cómputo nuevo de tienda y DC-01, los equipos de energía y recinto y el software bajo su alcance. El Cliente compra los equipos de red, de los centros de distribución y de terreno, y ejecuta las obras y el cableado. El proponente especifica, configura e integra los componentes y coordina las interfaces con los proveedores. Las capacidades iniciales son reservas de diseño sujetas a las pruebas declaradas.

El formulario usa los códigos de sitio de la Tabla 4.13 del Subdocumento 4. SIT-01 agrupa las 22 tiendas, CD-01 es el centro de distribución principal, CD-02 es el centro de distribución de Concepción, DC-01 es el centro de datos de casa matriz, DC-02 es su sala de respaldo, CLD-01 es la región Azure Chile Central, CLD-02 es la región Azure Brazil South y EXT-01 reúne a las contrapartes externas. Cada componente se atribuye a un ámbito, que puede ser Retail, filial emisora o frontera X-01, porque la separación entre la tienda y la Filial emisora alcanza a la infraestructura. La columna Origen indica el documento de la propuesta que selecciona el componente. Cuando el valor es T-11, este formulario precisa el producto o el nivel de una capacidad que los demás documentos dejaron remitida a él.

La Tabla~\mbox{\ref{tab:t11-1}} indica qué contiene cada apartado y con qué sección de la propuesta se corresponde.

\begin{longtable}[]{@{}
  >{\raggedright\arraybackslash}p{(\columnwidth - 4\tabcolsep) * \real{0.3000}}
  >{\raggedright\arraybackslash}p{(\columnwidth - 4\tabcolsep) * \real{0.4500}}
  >{\raggedright\arraybackslash}p{(\columnwidth - 4\tabcolsep) * \real{0.2500}}@{}}
\caption{Correspondencia entre el formulario y la propuesta.}\label{tab:t11-1}\\
\toprule\noalign{}
\textbf{Apartado} & \textbf{Contenido} & \textbf{Sección de la propuesta} \\
\midrule\noalign{}
\endfirsthead
\toprule\noalign{}
\textbf{Apartado} & \textbf{Contenido} & \textbf{Sección de la propuesta} \\
\midrule\noalign{}
\endhead
\bottomrule\noalign{}
\multicolumn{3}{@{}p{\columnwidth}@{}}{\rule{0pt}{3.2ex}\textbf{Fuente:} Elaboración propia.} \\
\endlastfoot
T-11.1 Servicios de nube contratados & Servicios de Azure por región, nivel, cantidad, ámbito y justificación. & SD-04, 4.1.1 y 4.2 \\ \addlinespace
T-11.2 Plataforma y software de la solución & Lenguajes, marcos, motores y software propio, con versión. & SD-04, 4.1.1 \\ \addlinespace
T-11.3 Seguridad, entrega y operación & Cadena de entrega, análisis de seguridad, monitoreo y respuesta. & SD-04, 4.2; SD-06, 6.2; Formulario T-10 \\ \addlinespace
T-11.4 Cómputo de tienda provisto por el proponente & Servidores locales, testigos, X-01 local y repuestos del suministro del proponente. & SD-04, 4.2.1 y Tabla 4.22 \\ \addlinespace
T-11.5 Red y equipos de sitio, hardware especificado por el proponente & Especificación de red de sitio, gabinete, energía, estaciones de trabajo, periféricos y terminales, que el Cliente adquiere. & SD-04, 4.2.1 y Tabla 4.22 \\ \addlinespace
T-11.6 Conectividad & Enlaces por sitio, conexión privada y respaldo. & SD-04, 4.2 y Tablas 4.13 y 4.22 \\ \addlinespace
T-11.7 Centros de datos & Centro primario, centro secundario, recinto técnico de DC-01 que habilita el proponente e instalaciones del Cliente. & SD-04, 4.3 y Tabla 4.22 \\ \addlinespace
T-11.8 Plataformas conservadas & Plataformas del Cliente que la solución integra y no oferta. & SD-04, 4.1 y 4.2.1 \\ \addlinespace
T-11.9 Componentes de las innovaciones & Equipos y plataformas de las cinco innovaciones. & SD-13 \\ \addlinespace
T-11.10 Licenciamiento & Modalidad, licencia, titular y vigencia por producto. & SD-04, 4.2.1 \\ \addlinespace
T-11.11 Ciclo de vida y soporte & Versión ofertada, fin de soporte publicado y plan de actualización. & SD-04, 4.1.1; SD-08 \\ \addlinespace
T-11.12 Portabilidad & Componentes portables y componentes propios del proveedor. & SD-04, 4.2; SD-08 \\ \addlinespace
T-11.13 Comprobaciones previas a la construcción & Verificaciones que fijan cantidades, modelos y niveles, con su hito. & SD-04, 4.2; SD-07 \\
\end{longtable}

Este formulario no contiene precios, tarifas ni montos que permitan inferir la Oferta Económica. Las cantidades que declara son las que alimentan esa oferta.

\subsection{T-11.1 Servicios de nube contratados}\label{t-11-1}

Microsoft Azure es el proveedor de nube (ADR-07). Los servicios se contratan bajo una suscripción del proyecto por ambiente y por ámbito, con Chile Central como región primaria y Brazil South como región de recuperación, ambas con tres zonas de disponibilidad. La Tabla~\mbox{\ref{tab:t11-2}} lista cada servicio con su nivel, su cantidad en el ambiente de Producción y la razón de su elección. Los cinco ambientes (Desarrollo, QA, Preproducción, Producción y Recuperación ante Desastres) usan las mismas clases de componentes. Preproducción y Recuperación tienen capacidades ajustadas a su finalidad y verificadas en la prueba de carga.

{\footnotesize
\begin{longtable}[]{@{}
  >{\raggedright\arraybackslash}p{(\columnwidth - 12\tabcolsep) * \real{0.0400}}
  >{\raggedright\arraybackslash}p{(\columnwidth - 12\tabcolsep) * \real{0.1200}}
  >{\raggedright\arraybackslash}p{(\columnwidth - 12\tabcolsep) * \real{0.1700}}
  >{\raggedright\arraybackslash}p{(\columnwidth - 12\tabcolsep) * \real{0.1600}}
  >{\raggedright\arraybackslash}p{(\columnwidth - 12\tabcolsep) * \real{0.1500}}
  >{\raggedright\arraybackslash}p{(\columnwidth - 12\tabcolsep) * \real{0.2400}}
  >{\raggedright\arraybackslash}p{(\columnwidth - 12\tabcolsep) * \real{0.1200}}@{}}
\caption{Servicios de Azure contratados.}\label{tab:t11-2}\\
\toprule\noalign{}
\textbf{N.°} & \textbf{Componente} & \textbf{Producto y nivel} & \textbf{Cantidad en Producción} & \textbf{Ubicación y ámbito} & \textbf{Justificación} & \textbf{Origen} \\
\midrule\noalign{}
\endfirsthead
\toprule\noalign{}
\textbf{N.°} & \textbf{Componente} & \textbf{Producto y nivel} & \textbf{Cantidad en Producción} & \textbf{Ubicación y ámbito} & \textbf{Justificación} & \textbf{Origen} \\
\midrule\noalign{}
\endhead
\bottomrule\noalign{}
\multicolumn{7}{@{}p{\columnwidth}@{}}{\rule{0pt}{3.2ex}\textbf{Fuente:} Elaboración propia a partir del Subdocumento 4 (Tablas 4.12, 4.9 y 4.10, y secciones 4.2 y 4.3), de su memoria de capacidad y del Formulario T-10.} \\
\endlastfoot
1.1 & Suscripciones & Suscripciones segregadas por ambiente y ámbito. & 15, tres ámbitos por cinco ambientes. & CLD-01 y CLD-02. Ámbito Retail, filial emisora y X-01. & Sin suscripciones compartidas entre ámbitos. & SD-04, 4.2.3 \\ \addlinespace
1.2 & Microservicios y adaptadores & AKS Premium con soporte de largo plazo propuesto y disponibilidad regional por comprobar. & Tres clústeres, 17 nodos Retail, 17 filial y 6 X-01, total 40 base, más 0 a 3 elásticos Retail. & CLD-01 y 40 nodos base de espera en CLD-02. Ámbito Retail, filial emisora y X-01. & Aislamiento de administración, capacidad y fallas por ámbito. & SD-04, Tablas 4.12 y 4.20 \\ \addlinespace
1.3 & Mensajería & Azure Service Bus Premium, colas, tópicos, sesiones y cola de mensajes no entregables. & Dos namespaces por región, uno por negocio. X-01 no utiliza bus propio. & CLD-01 y namespaces de recuperación en CLD-02. Ámbito Retail y filial emisora. & Hechos confirmados con deduplicación. Recuperación desde registros durables de productores y consumidores, sin asumir réplica de mensajes. & SD-04, 4.1.8, 4.3.2 y D-02 \\ \addlinespace
1.4 & Entrada API & API Management Premium clásico multirregión, sujeto a disponibilidad en ambas regiones. & Dos instancias, una por negocio, con gateways en Chile y Brasil. Unidades según capacidad y distribución zonal. & CLD-01 y CLD-02. Ámbito Retail y filial emisora. & Brasil atiende con la última configuración sincronizada sin publicar cambios durante la caída del plano primario. & SD-04, 4.1.1, 4.3.2 y D-03 \\ \addlinespace
1.5 & Datos operacionales y derivados & PostgreSQL Flexible Server, alta disponibilidad entre zonas para las autoridades y configuración por recurso. & 12 servidores lógicos, 6 Retail, 5 filial y 1 X-01. Reserva de 4.500 GiB, sin contar las esperas como autoridades nuevas. & CLD-01 y contrapartes asíncronas en CLD-02. Ámbito Retail, filial emisora y X-01. & Separa autoridades, analítica, vectores y metadatos. El retraso se mide y se prueba contra el RPO. & SD-04, Tabla 4.20 y Figura 4.22 \\ \addlinespace
1.6 & Caché & Azure Managed Redis, nivel y disponibilidad regional por comprobar. & Una instancia, reserva inicial de 16 GiB. & CLD-01. Ámbito Retail. & Solo información reconstruible, sin autoridad de negocio. & SD-04, Tablas 4.12 y 4.20 \\ \addlinespace
1.7 & Almacenamiento analítico & ADLS Gen2 por negocio, con redundancia de zona. & Dos lagos, 3 TiB Retail y 2 TiB filial como reservas iniciales. & CLD-01 y copia explícita en CLD-02. Ámbito Retail y filial emisora. & Historias y conjuntos autorizados separados por entidad. & SD-04, Tabla 4.20 \\ \addlinespace
1.8 & Analítica e informes & Data Factory, PostgreSQL analítico dedicado y Power BI. & Un entorno por negocio y cuatro VM de gateway, incluidas en las doce auxiliares. & CLD-01, con reconstrucción analítica y cuatro VM BI apagadas en CLD-02. Ámbito Retail y filial emisora. & Power BI consulta modelos analíticos propios, sin lectura directa de tablas operacionales. & SD-04, 4.1.10 y Tabla 4.20 \\ \addlinespace
1.9 & Evidencia y archivo & Blob de evidencia y archivo inmutable por clase de dato. & Evidencia filial de 5 TiB, archivo Retail de 25 TiB y filial de 12 TiB como reservas. & CLD-01 y copia explícita en CLD-02 con residencia aprobada. Ámbito Retail y filial emisora. & La evidencia contractual no se trata como proyección reconstruible. & SD-04, 4.1.10 y Tabla 4.20 \\ \addlinespace
1.10 & Identidad & Microsoft Entra ID, nivel P2, federado con el directorio del Cliente. & Identidades segregadas por ámbito; usuarios según el Caso (6.400 personas con acceso, 8.300 en período de mayor dotación). & Servicio global administrado desde CLD-01. Ámbito Retail, filial emisora y X-01. & Identidad centralizada con autenticación multifactor, acceso condicional y elevación de privilegios con vigencia limitada. & SD-04, Tabla 4.12; nivel T-11 \\ \addlinespace
1.11 & Secretos y claves & Key Vault por ámbito y ambiente. Managed HSM candidato para filial y X-01. & Bóvedas propias en cada región. HSM sujeto a censo de claves y viabilidad regional. & CLD-01 y CLD-02. Ámbito Retail, filial emisora y X-01. & Recuperación autorizada de claves, sin presuponer restauración directa de respaldos entre geografías. & SD-04, 4.2.6 y 4.3.2 \\ \addlinespace
1.12 & Observabilidad & OpenTelemetry Collector, Loki, Tempo, Prometheus y Grafana propios. Azure Monitor para recursos Azure. & Una pila por ámbito dentro de sus nodos base, incluida X-01 en sus seis nodos. Dos sondas externas ya contabilizadas. & CLD-01 y definición equivalente en CLD-02. Ámbito Retail, filial emisora y X-01. & Señales segregadas. Las sondas consultan salud mínima autorizada de X-01. & SD-04, 4.1.14 y D-06 \\ \addlinespace
1.13 & Registro de seguridad & Log Analytics y Sentinel por ámbito. & Tres espacios y configuraciones segregados por región de operación. & CLD-01 y CLD-02. Ámbito Retail, filial emisora y X-01. & Retención de seguridad de 12 meses en línea y 24 en archivo. No centraliza todos los registros de Loki. & SD-04, Tabla 4.20 \\ \addlinespace
1.14 & Borde público & Front Door Premium, WAF y Application Gateway por negocio. & Dos perfiles Front Door y un Application Gateway por negocio y región. & Borde global, CLD-01 y CLD-02. Ámbito Retail y filial emisora. & Entradas separadas. X-01 conserva su entrada privada de contrato. & SD-04, 4.2.3 y 4.2.4 \\ \addlinespace
1.15 & Control de red & Azure Firewall Premium. & Tres por región, uno por ámbito. & CLD-01 y CLD-02. Ámbito Retail, filial emisora y X-01. & Sin rutas directas entre negocios, cruce autorizado mediante X-01. & SD-04, Tabla 4.20 \\ \addlinespace
1.16 & Nombres y administración privada & DNS privado, DNS Private Resolver, Bastion y Entra Private Access. & Recursos propios por ámbito, cuatro VM de conectores Private Access entre las doce auxiliares. & CLD-01 y CLD-02. Ámbito Retail, filial emisora y X-01. & Acceso administrativo y resolución de nombres sin recursos compartidos entre negocios. & SD-04, 4.2.4 \\ \addlinespace
1.17 & Acceso privado a datos & Puntos de conexión privados de Azure Private Link. & Uno por servicio de datos en la subred de datos de cada radio. & CLD-01 y CLD-02. Ámbito Retail, filial emisora y X-01. & Ningún servicio de datos es alcanzable desde Internet. & SD-04, 4.2 \\ \addlinespace
1.18 & Redes virtuales & Hub y dos ramas por ámbito. & Nueve redes virtuales por región, tres por ámbito. & CLD-01 y CLD-02. Ámbito Retail, filial emisora y X-01. & Separación de redes y direcciones, sin peering entre negocios. & SD-04, 4.2.4 \\ \addlinespace
1.19 & Alternativa de enlace con DC-01 & ExpressRoute, alternativa adicional. & Sin circuito ni cantidad comprometida en la solución base. & DC-01 y CLD-01 si se adopta. Ámbito Por ámbito autorizado. & La VPN es la conexión base. ExpressRoute se evalúa contra volumen y criticidad. & SD-04, 4.2.4 y 4.2.8 \\ \addlinespace
1.20 & Conexión base de sitios & VPN Gateway con túneles IPsec segregados. & Dos gateways por región, uno por negocio. X-01 no recibe túneles de sitios sin necesidad acreditada. & CLD-01 y CLD-02. Ámbito Retail y filial emisora. & Dos accesos fijos y respaldo móvil por tienda, sin concentrador compartido. & SD-04, 4.2.4 y Tabla 4.20 \\ \addlinespace
1.21 & Respaldos & Bóveda de respaldo por ámbito y región, según cobertura del recurso. & Una por ámbito y región. & CLD-01 y CLD-02. Ámbito Retail, filial emisora y X-01. & Copia entre regiones explícita y restauración probada con sus claves. & SD-04, 4.2.6 \\ \addlinespace
1.22 & Imágenes de servicios & Container Registry Premium. & Tres registros, uno por ámbito, con copia identificada por digest a Brasil. & CLD-01 y CLD-02. Ámbito Retail, filial emisora y X-01. & Extracción de imágenes en contingencia sin depender de Chile. & SD-04, 4.2.5 \\ \addlinespace
1.23 & Cómputo predictivo & Grupo elástico AKS de Retail. & 0 a 3 nodos adicionales, fuera de los cuarenta base. & CLD-01. Ámbito Retail. & Cargas no críticas para la venta y alternativa manual. & SD-04, Tabla 4.20 \\ \addlinespace
1.24 & Módulo de seguridad de hardware & Managed HSM candidato para filial emisora y X-01. & Una capacidad dedicada por ámbito que lo requiera, sujeta a censo y disponibilidad. & Bóvedas y recuperación de claves propias en Chile y Brasil. & No se supone réplica automática entre regiones. Ensayo de descifrado con claves recuperadas. & SD-04, 4.2.6 y 4.3.2 \\ \addlinespace
1.25 & Distribución web & Application Gateway v2. & Uno por negocio y región. & Entrada desde Front Door, separada para Retail y filial emisora. & Entrega las solicitudes al gateway regional de API Management de su negocio. & SD-04, 4.2.4 y D-03 \\ \addlinespace
1.26 & Terminación de túneles & VPN Gateway, sin FortiGate-VM adicionales en la nube. & Incluido en la fila 1.20, sin doble cómputo. & Hubs propios de Retail y filial emisora. & Los FortiGate de sitio terminan sus túneles en VPN Gateway. Azure Firewall es el firewall de nube. & SD-04, 4.2.3 y 4.2.4 \\ \addlinespace
1.27 & Rutas entre redes & Rutas por ámbito y peering entre hub y ramas. & Según las nueve redes de la fila 1.18. No se compromete Route Server adicional. & Chile y Brasil, sin peering entre negocios. & El cruce autorizado se limita a los contratos de X-01. & SD-04, 4.2.4 \\ \addlinespace
1.28 & Administración privada & Bastion Premium y Entra Private Access. & Recursos por ámbito y cuatro conectores entre las doce VM auxiliares, sin contarlos otra vez. & CLD-01 y recuperación en CLD-02. & Acceso administrativo autorizado y segregado. & SD-04, 4.1.13 y 4.2.4 \\ \addlinespace
1.29 & Protección y salida & WAF, protección DDoS y DNS público, con salida por Azure Firewall. & Recursos separados por negocio, capacidad según prueba. & Borde global y hubs propios. & No se agrega un NAT Gateway compartido ni otra ruta entre negocios. & SD-04, 4.2.3 y 4.2.4 \\ \addlinespace
1.30 & Gobierno y operación & Azure Policy y gestión de costos por ámbito. & Quince suscripciones segregadas. Azure Arc solo si se justifica para un componente local. & Chile, Brasil y sitios bajo gestión autorizada. & Permisos, etiquetas y operación propios, sin una plataforma compartida de datos. & SD-04, 4.2.3 \\ \addlinespace
1.31 & Búsqueda y catálogo & OpenSearch global solo Retail. OpenMetadata con índice técnico por negocio. & Instancias lógicas aisladas dentro de la capacidad AKS reservada. & Grupos de estado de Retail y filial emisora. & La búsqueda técnica de catálogo de la filial no constituye búsqueda global de negocio. & SD-04, 4.1.10 y D-04 \\ \addlinespace
\end{longtable}
}

La Tabla T-11.2 distingue los tres ámbitos sin una plataforma de negocio compartida. Retail y la filial emisora tienen buses y entradas públicas propios. X-01 tiene clúster, bitácora, red, identidad y observabilidad propios, sin otro bus ni entrada pública. Las cantidades son reservas de diseño sujetas a las comprobaciones del SD-04.

La fila 1.6 declara Azure Managed Redis, como la Tabla 4.12 del Subdocumento 4, porque Microsoft dejó de admitir la creación de instancias de Azure Cache for Redis para clientes nuevos en abril de 2026 y fijó su retiro para el 30 de septiembre de 2028, dentro del período contractual (Microsoft, 2026a). Azure Managed Redis es el servicio que lo reemplaza y cumple la misma función. El Subdocumento 4 lo mantiene «por probar», y Valkey queda como alternativa por evaluar. Su creación en Chile Central se comprueba junto con los demás servicios (Tabla~\mbox{\ref{tab:t11-14}}).

Las unidades de API Management se dimensionan por gateway regional con distribución zonal y pruebas de carga. Las dos instancias multirregión no equivalen a dos instancias independientes por región. Los cuarenta nodos base, las doce VM auxiliares y las dos sondas suman 54 VM en Chile. En Brasil se mantienen cincuenta encendidas y cuatro de Power BI apagadas, según SD-04, 4.3.2.

\subsection{T-11.2 Plataforma y software de la solución}\label{t-11-2}

La solución nueva se construye con una pila única para los servicios y otra para las interfaces. La Tabla~\mbox{\ref{tab:t11-3}} lista el software de base y el software propio, con la versión ofertada y el lugar donde se ejecuta.

{\footnotesize
\begin{longtable}[]{@{}
  >{\raggedright\arraybackslash}p{(\columnwidth - 10\tabcolsep) * \real{0.0500}}
  >{\raggedright\arraybackslash}p{(\columnwidth - 10\tabcolsep) * \real{0.1400}}
  >{\raggedright\arraybackslash}p{(\columnwidth - 10\tabcolsep) * \real{0.2200}}
  >{\raggedright\arraybackslash}p{(\columnwidth - 10\tabcolsep) * \real{0.1700}}
  >{\raggedright\arraybackslash}p{(\columnwidth - 10\tabcolsep) * \real{0.2900}}
  >{\raggedright\arraybackslash}p{(\columnwidth - 10\tabcolsep) * \real{0.1300}}@{}}
\caption{Plataforma y software de la solución.}\label{tab:t11-3}\\
\toprule\noalign{}
\textbf{N.°} & \textbf{Componente} & \textbf{Producto y versión ofertada} & \textbf{Ubicación y ámbito} & \textbf{Justificación} & \textbf{Origen} \\
\midrule\noalign{}
\endfirsthead
\toprule\noalign{}
\textbf{N.°} & \textbf{Componente} & \textbf{Producto y versión ofertada} & \textbf{Ubicación y ámbito} & \textbf{Justificación} & \textbf{Origen} \\
\midrule\noalign{}
\endhead
\bottomrule\noalign{}
\multicolumn{6}{@{}p{\columnwidth}@{}}{\rule{0pt}{3.2ex}\textbf{Fuente:} Elaboración propia a partir del Subdocumento 4 (Tablas 4.2, 4.4, 4.6, 4.7 y 4.15).} \\
\endlastfoot
2.1 & Lenguaje de los servicios & Java 25 LTS, en una distribución OpenJDK con soporte. & AKS en CLD-01 y CLD-02. Ámbito Retail y filial emisora. & Versión de soporte prolongado vigente; una sola pila para APIs y lógica transaccional reduce la carga de operación y de competencias. & SD-04, Tabla 4.12 \\ \addlinespace
2.2 & Marco de los servicios & Spring Boot, línea 4, en su versión menor con soporte al inicio de la construcción. & AKS en CLD-01 y CLD-02. Ámbito Retail y filial emisora. & Marco compatible con Java 25 para contratos síncronos, eventos y acceso a datos. & SD-04, Tabla 4.12 \\ \addlinespace
2.3 & Lenguaje y biblioteca de interfaces & TypeScript y React, en sus versiones estables vigentes al inicio de la construcción. & Navegador, portales y vistas del punto de venta. Ámbito Retail y filial emisora. & Una sola base de componentes para portales internos, externos y caja. & SD-04, Tabla 4.12 \\ \addlinespace
2.4 & Orquestación de contenedores & Kubernetes, en una versión admitida por AKS. & CLD-01 y CLD-02. Ámbito Retail y filial emisora. & Aislamiento por entidad, escalado y recuperación con la misma definición. & SD-04, Tabla 4.12 \\ \addlinespace
2.5 & Motor relacional & PostgreSQL, en la versión mayor más reciente que el servicio administrado ofrezca en la región. & CLD-01 y CLD-02. Ámbito Retail y filial emisora. & Motor abierto, portable y con alta disponibilidad administrada. & SD-04, Tabla 4.12 \\ \addlinespace
2.6 & Mensajería de servicios & Clientes de Service Bus con contratos versionados. & Productores y consumidores de cada negocio. Ámbito Retail y filial emisora. & Entrega al menos una vez, orden cuando se exige, deduplicación y reconstrucción desde registros recuperados. & SD-04, 4.1.8 y 4.3.2 \\ \addlinespace
2.7 & Instrumentación & OpenTelemetry. & Todos los servicios y el nodo local. Ámbito Retail y filial emisora. & Estándar abierto de métricas, registros y trazas, independiente del destino. & SD-04, Tabla 4.12 \\ \addlinespace
2.8 & Infraestructura como código & Terraform. & Repositorios y canal de entrega segregados. Ámbito Retail, filial emisora y X-01. & Estado por ámbito y ambiente, con copia recuperable sin depender de Chile. & SD-04, 4.2.5 \\ \addlinespace
2.9 & Servicios de negocio & Trece servicios propios: R:M-01 a R:M-03, R:V-01 a R:V-04, R:CL-01, R:CL-02, F:C-01 a F:C-03 y X-01. & AKS en CLD-01. Ámbito Retail, filial emisora y X-01. & Desarrollo propio que implementa la partición de dominios (ADR-05). & SD-04, Tabla 4.3 \\ \addlinespace
2.10 & Adaptadores & Componentes propios en contenedores. & CLD-01, o DC-01 y CD-01 cuando la interfaz lo exija. Ámbito Retail y filial emisora. & Capa anticorrupción hacia cada plataforma conservada; la ruta se fija con el inventario de las catorce interfaces. & SD-04, Tablas 4.12 y 4.7 \\ \addlinespace
2.11 & Software del punto de venta & Aplicación propia de caja con modo local. & Terminales de caja de SIT-01. Ámbito Retail. & Sustituye el punto de venta de 2014 tienda por tienda (ADR-04). & SD-04, Tabla 4.22 \\ \addlinespace
2.12 & Software del nodo local & Caché de oferta (L-01), diario local (L-02) y sincronizador (L-03), desarrollo propio. & Nodos locales de SIT-01. Ámbito Retail. & Sostiene la venta y el cobro durante 24 horas sin enlace y concilia al reconectar. & SD-04, Tablas 4.12 y 4.20 \\ \addlinespace
2.13 & Registros durables de recuperación & Registro de salida por productor y registro de hechos aplicados por consumidor. & Datos del servicio propietario y sus copias recuperables. Ámbito Retail y filial emisora. & Conservación que cubre interrupción y reproceso, sin eliminar hechos solo por aceptación del bus. & SD-04, 4.3.2 y D-02 \\ \addlinespace
2.14 & Modelos y aplicaciones móviles & MLflow, Python y scikit-learn en Retail, y React Native para interfaces móviles nuevas. & Clúster Retail y dispositivos autorizados. & Versiones y empaquetado sujetos a las pruebas del SD-04. & SD-04, Tabla 4.12 \\ \addlinespace
2.15 & Observabilidad autogestionada & Collector, Loki, Tempo, Prometheus, Grafana y Blackbox Exporter. & Cargas propias en los tres clústeres y dos sondas externas ya contabilizadas. & X-01 utiliza sus seis nodos, con volúmenes y permisos propios. & SD-04, 4.1.14 y Tabla 4.20 \\ \addlinespace
\end{longtable}
}

La Tabla~\mbox{\ref{tab:t11-3}} muestra que los componentes de fabricante se concentran en la plataforma, mientras que los servicios de negocio, los adaptadores, el punto de venta y el nodo local son desarrollo propio cuya propiedad se transfiere al Cliente. El motor y el empaquetado del nodo local, junto con su sistema operativo, se fijan con el inventario de hardware y el ensayo de continuidad de 24 horas (Tabla~\mbox{\ref{tab:t11-14}}), que es la comprobación que la Tabla 4.12 exige antes de comprometerlos.

\subsection{T-11.3 Seguridad, entrega y operación}\label{t-11-3}

La cadena de entrega, el análisis de seguridad y el monitoreo usan las herramientas que declara el Formulario T-10, a las que este formulario agrega el producto de correlación de eventos y el de detección en puntos finales. La Tabla~\mbox{\ref{tab:t11-4}} las reúne con su función y su condición de soporte.

{\footnotesize
\begin{longtable}[]{@{}
  >{\raggedright\arraybackslash}p{(\columnwidth - 12\tabcolsep) * \real{0.0400}}
  >{\raggedright\arraybackslash}p{(\columnwidth - 12\tabcolsep) * \real{0.1400}}
  >{\raggedright\arraybackslash}p{(\columnwidth - 12\tabcolsep) * \real{0.1800}}
  >{\raggedright\arraybackslash}p{(\columnwidth - 12\tabcolsep) * \real{0.1200}}
  >{\raggedright\arraybackslash}p{(\columnwidth - 12\tabcolsep) * \real{0.1300}}
  >{\raggedright\arraybackslash}p{(\columnwidth - 12\tabcolsep) * \real{0.2600}}
  >{\raggedright\arraybackslash}p{(\columnwidth - 12\tabcolsep) * \real{0.1300}}@{}}
\caption{Herramientas de seguridad, entrega y operación.}\label{tab:t11-4}\\
\toprule\noalign{}
\textbf{N.°} & \textbf{Función} & \textbf{Producto} & \textbf{Ubicación} & \textbf{Condición de soporte} & \textbf{Justificación} & \textbf{Origen} \\
\midrule\noalign{}
\endfirsthead
\toprule\noalign{}
\textbf{N.°} & \textbf{Función} & \textbf{Producto} & \textbf{Ubicación} & \textbf{Condición de soporte} & \textbf{Justificación} & \textbf{Origen} \\
\midrule\noalign{}
\endhead
\bottomrule\noalign{}
\multicolumn{7}{@{}p{\columnwidth}@{}}{\rule{0pt}{3.2ex}\textbf{Fuente:} Elaboración propia a partir del Formulario T-10 y del Subdocumento 6.} \\
\endlastfoot
3.1 & Repositorio y canal de entrega & Azure Pipelines y repositorio versionado, con reutilización del canal del Cliente condicionada a revisión, firma y trazabilidad. & Agentes efímeros y registros por ámbito & Servicio contratado según selección final & Tres juegos de agentes e identidades, con permisos separados por ambiente. & SD-04, 4.2.3 y 4.2.5 \\ \addlinespace
3.2 & Análisis estático y calidad & SonarQube. & CLD-01 & Suscripción del fabricante & Bloquea el avance de una versión que no cumple los umbrales de calidad. & Formulario T-10 \\ \addlinespace
3.3 & Composición, imágenes e inventario de componentes & Trivy y Syft. & Canal de entrega & Comunidad, código abierto & Detecta vulnerabilidades en dependencias, imágenes e infraestructura como código, y genera el inventario de componentes. & Formulario T-10 \\ \addlinespace
3.4 & Firma de artefactos & Cosign. & Canal de entrega & Comunidad, código abierto & Solo se despliegan imágenes firmadas. & Formulario T-10 \\ \addlinespace
3.5 & Análisis dinámico & OWASP ZAP. & Ambiente de QA & Comunidad, código abierto & Prueba la aplicación en ejecución antes de Preproducción. & Formulario T-10 \\ \addlinespace
3.6 & Pruebas de contrato, extremo a extremo y carga & Pact, Playwright y k6. & QA y Preproducción & Comunidad, código abierto & Verifican contratos entre servicios, recorridos completos y la carga de 1,5 veces el máximo. & Formulario T-10 \\ \addlinespace
3.7 & Gestión del proyecto y documentación & Jira, Microsoft SharePoint y Slack. & Servicio en la nube & Suscripción del fabricante, contratada por el proponente & Herramientas de trabajo del proponente. No se transfieren al Cliente. & Formularios T-9 y T-10 \\ \addlinespace
3.8 & Correlación de seguridad & Sentinel sobre Log Analytics propios. & Un espacio por ámbito & Servicio administrado & No recibe automáticamente todos los registros de la pila operativa ni comparte datos entre negocios. & SD-04, 4.1.13 y 4.1.14 \\ \addlinespace
3.9 & Protección de cargas y de puntos finales & Microsoft Defender for Cloud y Microsoft Defender for Endpoint. & Nube, nodos locales y estaciones de trabajo de la Tabla~\mbox{\ref{tab:t11-6}} & Servicio administrado & Detección y respuesta en nube y en sitio con un solo proveedor de identidad y de registros. & T-11 \\ \addlinespace
3.10 & Configuración y despliegue local & Herramienta de configuración declarativa del nodo local, seleccionada entre las candidatas del Formulario T-10. & Nodos locales de SIT-01 & Según la herramienta seleccionada & Mantiene los 88 servidores de tienda segregados por negocio y con configuración reproducible y permite reconstruir uno desde el repositorio. & Formulario T-10 \\ \addlinespace
\end{longtable}
}

Los productos de las filas 3.2 a 3.6 son los de referencia del Formulario T-10. El Subdocumento 4 deja las herramientas de seguridad de la entrega como producto por definir, y el producto final se confirma en el diseño del canal de la Etapa 1. La Tabla~\mbox{\ref{tab:t11-4}} cubre los cuatro análisis que la matriz de cumplimiento compromete para el flujo de integración continua, que son el estático, el dinámico, el de composición y el de imágenes. Las filas 3.8 y 3.9 siguen el criterio de ADR-07, que privilegia la continuidad de un solo proveedor para identidad, datos y operación. La herramienta de la fila 3.10 se compromete con el diseño del nodo local, como establece el Formulario T-10.

\subsection{T-11.4 Cómputo de tienda provisto por el proponente}\label{t-11-4}

Only Simple Solutions provee los servidores y testigos de tienda conforme a la Tabla 4.22 del SD-04. El Cliente aporta espacio, acceso y permisos. La Tabla~\mbox{\ref{tab:t11-5}} distingue el parque instalado de sus repuestos. Se mantienen modelos de referencia por confirmar mediante la prueba de continuidad.

\begin{longtable}[]{@{}
  >{\raggedright\arraybackslash}p{(\columnwidth - 4\tabcolsep) * \real{0.2000}}
  >{\raggedright\arraybackslash}p{(\columnwidth - 4\tabcolsep) * \real{0.5000}}
  >{\raggedright\arraybackslash}p{(\columnwidth - 4\tabcolsep) * \real{0.3000}}@{}}
\caption{Cómputo de tienda provisto por Only Simple Solutions.}\label{tab:t11-5}\\
\toprule\noalign{}
\textbf{Atributo} & \textbf{4.1 Nodo local de tienda} & \textbf{4.2 Repuestos de nodo local} \\
\midrule\noalign{}
\endfirsthead
\toprule\noalign{}
\textbf{Atributo} & \textbf{4.1 Nodo local de tienda} & \textbf{4.2 Repuestos de nodo local} \\
\midrule\noalign{}
\endhead
\bottomrule\noalign{}
\multicolumn{3}{@{}p{\columnwidth}@{}}{\rule{0pt}{3.2ex}\textbf{Fuente:} Elaboración propia a partir del Subdocumento 4 (Tabla 4.22 y sección 4.2.1), del Subdocumento 6 y del Caso 09 (capítulo 15, p. 34).} \\
\endlastfoot
Quién especifica y quién adquiere & Only Simple Solutions especifica, provee, instala y mantiene el cómputo local y el software, conforme a SD-04, Tabla 4.22. & Only Simple Solutions provee y mantiene el stock compatible. \\ \addlinespace
Marca y modelo de referencia & Servidor compacto de borde. Marca y modelo de referencia por fijar con el ensayo de continuidad de 24 horas, en los meses 6 y 7. & El mismo modelo de la fila 4.1, con configuración precargada. \\ \addlinespace
Cantidad & 88, cuatro por tienda, 44 Retail y 44 filial emisora. & Diez, cinco por cada perfil de 44 servidores. \\ \addlinespace
Características mínimas & Retail, 8 núcleos y 32 GiB por servidor. Filial emisora, 12 núcleos y 48 GiB. Dos SSD de 960 GB en espejo, cifrado, endurecimiento y gestión remota. & Mismas características del perfil reemplazado. \\ \addlinespace
Consumo y accesorios & Consumo por fijar con el modelo de referencia, en los meses 6 y 7. Incluye rieles de montaje en gabinete y cables de alimentación y de red. & Igual que la fila 4.1. \\ \addlinespace
Margen de crecimiento & El dimensionamiento cubre la proyección a tres años del Caso. El porcentaje de margen sobre esa carga se declara con el modelo, en los meses 6 y 7. & No aplica \\ \addlinespace
Condiciones de uso y ciclo de vida & Uso en interior, dentro del gabinete de la tienda. Ciclo de vida de 56 meses, con garantía y soporte del fabricante por todo el contrato. & Igual que la fila 4.1. \\ \addlinespace
Redundancia & Un par por negocio en cada tienda, sin compartir servidor físico entre negocios. Cada servidor tolera la falla de un disco. & No integra el parque en producción. \\ \addlinespace
Ubicación & Gabinete de cada tienda de SIT-01. & Bodega en Chile, con despacho a cada tienda. \\ \addlinespace
Soporte & Provisión y mantenimiento del proponente con garantía del fabricante. Atención según SLA contractual aplicable. & Reemplazo compatible y configuración precargada, según plan logístico. \\ \addlinespace
Justificación & La tienda vende y cobra sin enlace durante 24 horas (ADR-04), y la falla de un nodo no detiene la caja. & Existencia de reemplazo del 10 \% del parque de componentes críticos, disponible en Chile (RT-08.4). \\
\end{longtable}

Los 88 servidores resultan de 22 tiendas por cuatro servidores y constituyen una decisión de diseño, no una cantidad impuesta por el Caso. Los diez repuestos resultan de dos perfiles de 44 equipos, con cinco repuestos por perfil. Con 25 tiendas, el parque sería de cien servidores y se mantendrían cinco repuestos por perfil de cincuenta. Los testigos y demás componentes críticos requieren su propio stock conforme al inventario final. Los modelos, la potencia y la capacidad se confirman mediante inventario y ensayo de continuidad. La pérdida de un servidor debe permitir operar con el otro del mismo negocio. Los dispositivos de red, los pares de captura de los centros y las obras se asignan al Cliente en el apartado siguiente.

Los servidores se complementan con 44 testigos, uno por par y negocio, y 44 instancias locales de X-01. Cada instancia reserva 2 vCPU y 4 GiB dentro de un servidor de la filial emisora, con bitácora propia de 20 GB, sin otro servidor físico. Los tres racks, las dos UPS y las seis PDU por tienda son suministro del proponente. Los 24 routers móviles, uno por tienda y centro de distribución, son adquisición del Cliente. Estas cantidades se cuentan una sola vez.

\subsection{T-11.5 Red y equipos de sitio, hardware especificado por el proponente}\label{t-11-5}

El Caso establece que el hardware de terreno, el equipamiento de red y las obras los adquiere y ejecuta el Cliente, y que el proponente especifica qué comprar, cuánto y con qué características (Multitiendas Ancoa S.A., 2026b, cap. 11, p. 26). La Tabla~\mbox{\ref{tab:t11-6}} entrega esa especificación. Only Simple Solutions configura y certifica cada equipo antes de la ola de despliegue de su tienda, y el Cliente lo adquiere e instala. Donde el Subdocumento 4 no da una cantidad, la tabla escribe «por definir» con el hito del levantamiento. Los modelos de referencia se fijan en el paquete 1.4.8 de la EDT, antes del hito H3 (mes 6).

{\footnotesize
\begin{longtable}[]{@{}
  >{\raggedright\arraybackslash}p{(\columnwidth - 12\tabcolsep) * \real{0.0400}}
  >{\raggedright\arraybackslash}p{(\columnwidth - 12\tabcolsep) * \real{0.1300}}
  >{\raggedright\arraybackslash}p{(\columnwidth - 12\tabcolsep) * \real{0.1400}}
  >{\raggedright\arraybackslash}p{(\columnwidth - 12\tabcolsep) * \real{0.1300}}
  >{\raggedright\arraybackslash}p{(\columnwidth - 12\tabcolsep) * \real{0.2600}}
  >{\raggedright\arraybackslash}p{(\columnwidth - 12\tabcolsep) * \real{0.1100}}
  >{\raggedright\arraybackslash}p{(\columnwidth - 12\tabcolsep) * \real{0.1900}}@{}}
\caption{Especificación de hardware y obras que adquiere el Cliente.}\label{tab:t11-6}\\
\toprule\noalign{}
\textbf{N.°} & \textbf{Componente} & \textbf{Referencia de marca} & \textbf{Cantidad} & \textbf{Características mínimas} & \textbf{Ubicación} & \textbf{Justificación} \\
\midrule\noalign{}
\endfirsthead
\toprule\noalign{}
\textbf{N.°} & \textbf{Componente} & \textbf{Referencia de marca} & \textbf{Cantidad} & \textbf{Características mínimas} & \textbf{Ubicación} & \textbf{Justificación} \\
\midrule\noalign{}
\endhead
\bottomrule\noalign{}
\multicolumn{7}{@{}p{\columnwidth}@{}}{\rule{0pt}{3.2ex}\textbf{Fuente:} Elaboración propia a partir del Subdocumento 4 (Tabla 4.22, sección 4.2.1 y Figura 4.19) y del Caso 09 (capítulo 6, p. 14, capítulo 11, p. 26, capítulo 14, pp. 31-32 y capítulo 15, pp. 33-36).} \\
\endlastfoot
5.1 & Conmutadores de tienda & Fabricante con soporte en Chile y características equivalentes. & 88 posiciones, cuatro por tienda. Compra neta según reutilización aceptada. & Dos por negocio, segmentación y conmutación sin caída por pérdida de uno. & 22 tiendas & Patrón de la Figura 4.19 del SD-04. \\ \addlinespace
5.2 & Firewalls de tienda & Fabricante con soporte en Chile y características equivalentes. & 88 posiciones, cuatro por tienda. Compra neta según inventario. & Un par por negocio, SD-WAN, túneles segregados y conmutación automática. & 22 tiendas & Separación Retail y filial, sin firewall compartido. \\ \addlinespace
5.3 & Racks de tienda & Tres racks de 24U por tienda. & 66 en total, incluidos en el suministro del proponente, sin compra adicional del Cliente. & Dos racks de servidores y uno de comunicaciones, 31 de 72U iniciales. & SIT-01 & SD-04, 4.2.1 y Tabla 4.22. \\ \addlinespace
5.4 & Energía de tienda & Dos UPS de 6 kW y seis PDU por tienda. & 44 UPS y 132 PDU, del suministro del proponente. & Cada UPS sostiene la carga completa durante treinta minutos. Clima N+1 y transferencia según SD-04. & SIT-01 & El Cliente ejecuta obras eléctricas y canalizaciones. No duplicar compra de equipos. \\ \addlinespace
5.5 & Puntos de acceso Wi-Fi & Modelo por fijar antes del hito H3 (mes 6). & Dos tipos por tienda, uno corporativo y uno aislado para clientes. Cantidad de puntos por tienda por definir en el levantamiento, meses 3 a 6. & El corporativo solo admite equipos con certificado de identidad y atiende la red de móviles. El de clientes queda aislado de las demás redes. & SIT-01 & Mantiene el Wi-Fi de clientes fuera de la red de cajas y protege los datos de pago y de crédito. \\ \addlinespace
5.6 & Cámaras y grabador de video & Se conservan los existentes. & Según el inventario por tienda. & Red propia de video y grabador con 30 días de imágenes. Si el inventario muestra que no cumple, se agrega su especificación. & SIT-01 & Mantiene la videovigilancia en su propia red. \\ \addlinespace
5.7 & Terminales de caja & Se conservan los existentes. Marca y modelo por inventariar. & 380 líneas de caja, 430 en la proyección a tres años. & Compatibilidad con el software del punto de venta nuevo, verificada en el prototipo. & SIT-01 & El punto de venta se sustituye por software, tienda por tienda, sin cambiar el parque de cajas. \\ \addlinespace
5.8 & Periféricos de caja & Se conservan los existentes. Marca y modelo por inventariar. & Un conjunto por línea de caja, 380. & Impresora fiscal, lector de código y terminal de pago, con controladores verificados en el prototipo. & SIT-01 & La Tabla 4.22 los mantiene a cargo del Cliente. \\ \addlinespace
5.9 & Terminales móviles de sala & Se conservan los existentes. Marca y modelo por inventariar. & 640, compartidas por 3.820 vendedores. & Sistema operativo vigente o una de sus dos versiones anteriores, gestión remota de flota y operación sin conexión para los perfiles internos. & SIT-01 & El Caso no pide adquirirlas y sí especificar la dotación. La dotación requerida se calcula con el uso por turno medido en las tiendas piloto. \\ \addlinespace
5.10 & Impresoras de etiqueta & Se conservan las existentes. Marca y modelo por inventariar. & Al menos una por tienda, 22. & Compatibilidad con la confirmación de actualización de precio. & Bodega de cada tienda de SIT-01 & Sostienen el cambio de precio físico mientras no haya etiqueta electrónica, que el Caso excluye de la instalación. \\ \addlinespace
5.11 & Red de centros & Modelos por levantamiento. & Según inventario, con un router móvil por centro. & Red propia de Retail, túneles y captura local. & CD-01 y CD-02 & Cliente adquiere e instala, proponente especifica, configura e integra. \\ \addlinespace
5.12 & Captura local de centros & Par local por centro, distinto del WMS conservado. & Cuatro servidores y cuatro racks de 24U en total. & Un par de captura y dos racks por centro, sin instalar otro WMS en Concepción. & CD-01 y CD-02 & Compra e instalación del Cliente, integración del proponente. \\ \addlinespace
5.13 & Estaciones de trabajo de operación y administración & Modelo de referencia por fijar antes del hito H3 (mes 6). & Por definir con el dimensionamiento del equipo de operación, hito H2 (mes 4). & Dos monitores por estación, gestión central, cifrado de disco, control de dispositivos extraíbles y detección y respuesta en el punto final. Ergonomía conforme a la NCh 2527 y certificación de eficiencia energética (RT-08.07 a RT-08.09). & Puestos de operación del proponente y de la Contraparte Técnica & La administración de la plataforma requiere puestos con la protección de la fila 3.9. \\ \addlinespace
5.14 & Obras de sitio & No aplica. & Según el inventario por tienda. & Canalizaciones, circuitos eléctricos independientes y cableado estructurado certificado para el gabinete. & SIT-01, CD-01 y CD-02 & El Caso asigna la ejecución al Cliente y la especificación al proponente. \\
\end{longtable}
}

La Tabla~\mbox{\ref{tab:t11-6}} distingue los equipos que adquiere el Cliente de los racks y UPS que aporta el proponente. El objetivo de red es de cuatro switches y cuatro firewalls por tienda. La compra neta descuenta equipos existentes solo si su capacidad, soporte y aislamiento se aceptan técnicamente. Las filas de equipos de energía y recinto se mantienen para documentar su interfaz con las obras del Cliente, sin duplicar su adquisición.

El Cliente conserva las garantías y la reparación física de los equipos que adquiere. Only Simple Solutions mantiene las configuraciones y el software bajo su alcance. El stock de reemplazo de equipos críticos se calcula por perfil y se redondea hacia arriba. Los 88 switches y 88 firewalls del parque objetivo requieren al menos nueve repuestos de cada tipo si son homogéneos. Las capacidades y el reparto por modelo se confirman en el inventario. Los repuestos de servidores figuran en la Tabla~\mbox{\ref{tab:t11-5}}.

\subsection{T-11.6 Conectividad}\label{t-11-6}

La conectividad une los sitios del Cliente con la nube por rutas privadas o cifradas. La Tabla~\mbox{\ref{tab:t11-7}} lista los enlaces por sitio y distingue los existentes de los que la propuesta requiere.

{\footnotesize
\begin{longtable}[]{@{}
  >{\raggedright\arraybackslash}p{(\columnwidth - 10\tabcolsep) * \real{0.0500}}
  >{\raggedright\arraybackslash}p{(\columnwidth - 10\tabcolsep) * \real{0.1500}}
  >{\raggedright\arraybackslash}p{(\columnwidth - 10\tabcolsep) * \real{0.2000}}
  >{\raggedright\arraybackslash}p{(\columnwidth - 10\tabcolsep) * \real{0.2500}}
  >{\raggedright\arraybackslash}p{(\columnwidth - 10\tabcolsep) * \real{0.1200}}
  >{\raggedright\arraybackslash}p{(\columnwidth - 10\tabcolsep) * \real{0.2300}}@{}}
\caption{Enlaces por sitio.}\label{tab:t11-7}\\
\toprule\noalign{}
\textbf{N.°} & \textbf{Sitio} & \textbf{Enlace existente} & \textbf{Enlace requerido por la solución} & \textbf{Cantidad nueva} & \textbf{Provee} \\
\midrule\noalign{}
\endfirsthead
\toprule\noalign{}
\textbf{N.°} & \textbf{Sitio} & \textbf{Enlace existente} & \textbf{Enlace requerido por la solución} & \textbf{Cantidad nueva} & \textbf{Provee} \\
\midrule\noalign{}
\endhead
\bottomrule\noalign{}
\multicolumn{6}{@{}p{\columnwidth}@{}}{\rule{0pt}{3.2ex}\textbf{Fuente:} Elaboración propia a partir del Subdocumento 4 (sección 4.2 y Tablas 4.13, 4.20 y 4.22) y del Caso 09 (capítulo 6, p. 14; capítulo 15, p. 34).} \\
\endlastfoot
6.1 & 8 tiendas en centro comercial & Infraestructura del centro y respaldo propio. & Dos accesos fijos de operadores y caminos distintos, mínimo 10 Mbps por acceso, más router móvil de 5 Mbps y túneles por negocio. & Según levantamiento, sin presuponer enlaces adicionales donde ya cumplen. & Cliente contrata, operadores proveen, proponente configura. \\ \addlinespace
6.2 & 6 tiendas en centro comercial & Infraestructura del centro sin respaldo propio. & Dos accesos fijos de operadores y caminos distintos, mínimo 10 Mbps por acceso, más router móvil de 5 Mbps y túneles por negocio. & Según levantamiento, sin presuponer enlaces adicionales donde ya cumplen. & Cliente contrata, operadores proveen, proponente configura. \\ \addlinespace
6.3 & 8 tiendas a la calle & Acceso propio, respaldo de cada tienda por acreditar. & Dos accesos fijos de operadores y caminos distintos, mínimo 10 Mbps por acceso, más router móvil de 5 Mbps y túneles por negocio. & Según levantamiento, sin presuponer enlaces adicionales donde ya cumplen. & Cliente contrata, operadores proveen, proponente configura. \\ \addlinespace
6.4 & CD-01 & Dos accesos de proveedores distintos. & Túneles de Retail sobre ambos accesos y respaldo móvil. Reserva de 2 Mbps para la solución. & Accesos fijos según inventario, router en 5.3. & Cliente y operadores, configuración del proponente. \\ \addlinespace
6.5 & CD-02 & Acceso único. & Segundo acceso independiente y respaldo móvil, reserva de 2 Mbps para la solución. & Un segundo acceso propuesto y router en 5.3. & Cliente y operadores, configuración del proponente. \\ \addlinespace
6.6 & DC-01 & Accesos y rutas por levantar. & Dos accesos con VPN por ámbito, dimensionados para 110 Mbps de migración. ExpressRoute es alternativa adicional. & Según levantamiento. Sin circuito ExpressRoute comprometido. & Cliente contrata accesos, proponente configura VPN. \\ \addlinespace
6.7 & DC-01 y DC-02 & Enlace no acreditado. & Conexión propuesta según función comprobada de DC-02. & Por acreditar y dimensionar. & Cliente contrata, proponente especifica e integra. \\ \addlinespace
6.8 & Oficinas & Acceso existente. & Túnel cifrado o acceso por red definida por software. & 0 & Cliente, enlaces; proponente, configuración \\ \addlinespace
6.9 & DC-01 y recuperación & Ruta de contingencia por acreditar. & VPN hacia el gateway regional del negocio recuperado. No se requiere un hub compartido. & Túneles por ámbito conforme a rutas autorizadas. & Proponente configura, Cliente contrata los accesos. \\ \addlinespace
\end{longtable}
}

La Tabla T-11.7 no fija una compra de catorce enlaces de respaldo. El inventario determina cuántos accesos faltan para disponer de dos caminos fijos por tienda y de proveedores independientes. Los 24 routers de la fila 5.3 se cuentan una sola vez. La VPN es la base de conexión y ExpressRoute no integra la cantidad comprometida. DC-02 no sustituye la recuperación regional.

\subsection{T-11.7 Centros de datos}\label{t-11-7}

La región primaria y la región de recuperación de Azure cumplen el papel de centro de datos principal y secundario de la solución nueva. La Tabla~\mbox{\ref{tab:t11-8}} compara sus especificaciones y agrega las instalaciones del Cliente, que permanecen en convivencia.

\begin{longtable}[]{@{}
  >{\raggedright\arraybackslash}p{(\columnwidth - 6\tabcolsep) * \real{0.1600}}
  >{\raggedright\arraybackslash}p{(\columnwidth - 6\tabcolsep) * \real{0.2700}}
  >{\raggedright\arraybackslash}p{(\columnwidth - 6\tabcolsep) * \real{0.2700}}
  >{\raggedright\arraybackslash}p{(\columnwidth - 6\tabcolsep) * \real{0.3000}}@{}}
\caption{Especificaciones de los centros de datos.}\label{tab:t11-8}\\
\toprule\noalign{}
\textbf{Atributo} & \textbf{Primario, CLD-01} & \textbf{Secundario, CLD-02} & \textbf{Instalaciones del Cliente, DC-01 y DC-02} \\
\midrule\noalign{}
\endfirsthead
\toprule\noalign{}
\textbf{Atributo} & \textbf{Primario, CLD-01} & \textbf{Secundario, CLD-02} & \textbf{Instalaciones del Cliente, DC-01 y DC-02} \\
\midrule\noalign{}
\endhead
\bottomrule\noalign{}
\multicolumn{4}{@{}p{\columnwidth}@{}}{\rule{0pt}{3.2ex}\textbf{Fuente:} Elaboración propia a partir del Subdocumento 4 (sección 4.3 y Tablas 4.13 y 4.22) y del Caso 09 (capítulo 6, p. 14).} \\
\endlastfoot
Identificación & Región Azure Chile Central (\texttt{chilecentral}). & Región Azure Brazil South (\texttt{brazilsouth}). & Centro de datos de casa matriz de 140 m² y sala de respaldo. \\ \addlinespace
Ubicación & Santiago, Chile. & São Paulo, Brasil. & Misma comuna ambos. \\ \addlinespace
Zonas de disponibilidad & Tres. & Tres. & No aplica. \\ \addlinespace
Función & Producción de los tres ámbitos. & Recuperación activo-pasivo templada, cincuenta VM encendidas y cuatro BI apagadas. & Convivencia e integración local. DC-02 no recupera las funciones nuevas. \\ \addlinespace
Distancia y amenazas comunes & Países distintos. La distancia exacta se declara con su fuente antes de aprobar el diseño detallado (RT-07.02). & Países distintos. La distancia exacta se declara con su fuente antes de aprobar el diseño detallado (RT-07.02). & Comparten comuna, por lo que DC-02 no basta como única recuperación. \\ \addlinespace
Componentes & Filas 1.1 a 1.23 y software propio. & Gateway regional APIM, 12 PostgreSQL lógicos, namespaces Service Bus con reconstrucción durable, objetos y claves recuperables, 40 nodos AKS y auxiliares. & Cómputo nuevo de integración en DC-01 y sistemas conservados con sitio supuesto o por acreditar. \\ \addlinespace
Replicación & Redundancia de zona según recurso. & Réplica de bases y copia explícita de objetos, conservación de registros de recuperación de mensajes. & Enlace y función de DC-02 por acreditar. \\ \addlinespace
Objetivos de recuperación & No aplica. & RTO no superior a 4 horas y RPO no superior a 15 minutos para servicios críticos. & No sostienen funciones nuevas como única ubicación. \\ \addlinespace
Ensayo & Falla de zona. & Conmutación y retorno dos veces al año, con carga y reproceso sin efectos duplicados. & Conmutación de accesos VPN y prueba de contingencia por plataforma. \\ \addlinespace
Energía, climatización y seguridad física & Infraestructura física del proveedor de nube. & Infraestructura física del proveedor de nube. & Cliente ejecuta obras. Only Simple Solutions habilita el recinto técnico y su integración según Tabla 4.22 del SD-04. \\ \addlinespace
\end{longtable}

La Tabla T-11.8 corresponde a la arquitectura base Chile y Brasil. El proponente incorpora dos pares nuevos de integración en DC-01 y no traslada allí las bases financieras nuevas. Crédito 2011 tiene ubicación por acreditar. La residencia se aprueba por ámbito antes del cierre del diseño detallado. Si Brasil no se autoriza, el proponente habilita y ensaya un sitio chileno independiente antes de la producción de la etapa correspondiente, con RTO máximo de cuatro horas y RPO máximo de quince minutos. Esa alternativa conserva separación y reconstrucción durable, sin depender de una futura región de Azure.

El proponente habilita técnicamente DC-01 y provee dos pares nuevos de integración, uno por negocio. El Cliente aporta el espacio, adquiere e instala la red y ejecuta obras y cableado, conforme a la Tabla 4.22 del SD-04. La Tabla~\mbox{\ref{tab:t11-dc01}} conserva el detalle del recinto sin atribuir al proponente las obras del Cliente. La carga y las brechas se confirman con el inventario y el informe interno de 2024.

{\footnotesize
\begin{longtable}[]{@{}
  >{\raggedright\arraybackslash}p{(\columnwidth - 10\tabcolsep) * \real{0.0500}}
  >{\raggedright\arraybackslash}p{(\columnwidth - 10\tabcolsep) * \real{0.1500}}
  >{\raggedright\arraybackslash}p{(\columnwidth - 10\tabcolsep) * \real{0.3400}}
  >{\raggedright\arraybackslash}p{(\columnwidth - 10\tabcolsep) * \real{0.2000}}
  >{\raggedright\arraybackslash}p{(\columnwidth - 10\tabcolsep) * \real{0.1500}}
  >{\raggedright\arraybackslash}p{(\columnwidth - 10\tabcolsep) * \real{0.1100}}@{}}
\caption{Equipos y obras de DC-01 y reparto de suministro.}\label{tab:t11-dc01}\\
\toprule\noalign{}
\textbf{N.°} & \textbf{Componente} & \textbf{Características mínimas} & \textbf{Cantidad} & \textbf{Referencia de marca} & \textbf{Requisito} \\
\midrule\noalign{}
\endfirsthead
\toprule\noalign{}
\textbf{N.°} & \textbf{Componente} & \textbf{Características mínimas} & \textbf{Cantidad} & \textbf{Referencia de marca} & \textbf{Requisito} \\
\midrule\noalign{}
\endhead
\bottomrule\noalign{}
\multicolumn{6}{@{}p{\columnwidth}@{}}{\rule{0pt}{3.2ex}\textbf{Fuente:} Elaboración propia a partir del Subdocumento 4 (Tabla 4.22 y sección 4.3.1) y de las Bases Técnicas Transversales (capítulo 6).} \\
\endlastfoot
7.1 & Firewall del recinto & Un par de firewall por negocio, sin par compartido. & Cuatro equipos, dos por negocio, sujetos a inventario. & Cliente adquiere e instala, proponente especifica y configura. & RT-08.03 y SD-04, Tabla 4.22 \\ \addlinespace
7.2 & Switches del recinto & Redes y administración propias por negocio. & Un par por negocio, cuatro equipos como patrón de diseño. & Cliente adquiere e instala, proponente especifica y configura. & RT-08.03 y SD-04, Figura 4.25 \\ \addlinespace
7.3 & Racks & Cuatro racks de 42U separados para servidores, comunicaciones y respaldo. & Cuatro, con reserva inicial de 80 de 168U. & Provisión del proponente, modelo según inventario. & RT-06.05 \\ \addlinespace
7.4 & Almacenamiento en red y respaldo local & Una zona Retail y otra del Emisor, cada una con su almacenamiento en red y su respaldo local, redundantes y con tolerancia declarada a la falla de discos. & 2 zonas. Capacidad por definir con el inventario de plataformas, hito H2 (mes 4). & Modelo por fijar en el hito H3 (mes 6). & RT-08.02 y línea roja del Caso \\ \addlinespace
7.5 & Cinta en custodia & Copia fuera de línea de DC-01, parte de la política de respaldo 3-2-1-1-0. & Por definir con el volumen de respaldo, hito H2 (mes 4). & Modelo por fijar en el hito H3 (mes 6). & Subdocumento 4, sección 4.2 \\ \addlinespace
7.6 & Energía ininterrumpida & Autonomía de al menos 30 minutos a plena carga. & Por definir con la carga eléctrica proyectada, hito H2 (mes 4). & Modelo por fijar en el hito H3 (mes 6). & RT-06.07 \\ \addlinespace
7.7 & Generación autónoma & Al menos 24 horas continuas, con combustible y contrato de reabastecimiento. & Por definir con la carga eléctrica proyectada, hito H2 (mes 4). & Modelo por fijar en el hito H3 (mes 6). & RT-06.08 \\ \addlinespace
7.8 & Climatización de precisión & Equipo de reserva (N+1). & Por definir con la carga térmica, hito H2 (mes 4). & Modelo por fijar en el hito H3 (mes 6). & RT-06.13 \\ \addlinespace
7.9 & Detección y extinción de incendio & Detección temprana por aspiración de aire y extinción automática por agente limpio. & Una instalación para el recinto. & Modelo por fijar en el hito H3 (mes 6). & RT-06.16 y RT-06.17 \\ \addlinespace
7.10 & Control de acceso & Biométrico. & Por definir en el levantamiento, hito H2 (mes 4). & Modelo por fijar en el hito H3 (mes 6). & RT-06.20 \\ \addlinespace
7.11 & Videovigilancia & Imágenes con 30 días de conservación. & Por definir en el levantamiento, hito H2 (mes 4). & Modelo por fijar en el hito H3 (mes 6). & RT-06.24 \\ \addlinespace
7.12 & Conectividad y obras & Conectividad segregada y canalizaciones según el plano del recinto. & Según proyecto por sitio. & Cliente ejecuta obras y cableado. Proponente especifica, coordina, configura e integra. & RT-06.03, RT-06.04 y RT-06.33 \\ \addlinespace
7.13 & Obra civil de separación y espacio & Separación del espacio de operación y de la sala de equipos. La obra y el espacio los aporta el Cliente y el proponente los especifica y coordina. & Una obra, según el informe interno de 2024. & No aplica. & RT-06.06 y RT-06.30 \\
\end{longtable}
}

La tipología que se adopta en DC-01 es la de sala técnica principal, a la que se aplican íntegramente RT-06.01 a RT-06.24, con una disponibilidad mensual de 99,95 \% para la energía, la climatización y la red del recinto. En DC-02 y en CD-01 se propone la de sala técnica secundaria o de sitio, y en cada tienda y en CD-02 la de gabinete o borde operacional. Esas dos últimas asignaciones se confirman con el levantamiento de la Etapa 1 (Tabla~\mbox{\ref{tab:t11-14}}). La solución nueva no depende de DC-01 ni de DC-02 como única ubicación. Los equipos de la tabla sostienen las plataformas conservadas y la conexión privada.

\subsection{T-11.8 Plataformas conservadas}\label{t-11-8}

Las plataformas actuales del Cliente no forman parte de lo ofertado. La Tabla~\mbox{\ref{tab:t11-9}} las lista porque la solución se integra con ellas y porque su hardware y sus licencias no se atribuyen a los servicios nuevos.

\begin{longtable}[]{@{}
  >{\raggedright\arraybackslash}p{(\columnwidth - 6\tabcolsep) * \real{0.2500}}
  >{\raggedright\arraybackslash}p{(\columnwidth - 6\tabcolsep) * \real{0.2500}}
  >{\raggedright\arraybackslash}p{(\columnwidth - 6\tabcolsep) * \real{0.2000}}
  >{\raggedright\arraybackslash}p{(\columnwidth - 6\tabcolsep) * \real{0.3000}}@{}}
\caption{Plataformas del Cliente y su tratamiento.}\label{tab:t11-9}\\
\toprule\noalign{}
\textbf{Plataforma} & \textbf{Destino} & \textbf{Alojamiento supuesto} & \textbf{Aporte del proponente} \\
\midrule\noalign{}
\endfirsthead
\toprule\noalign{}
\textbf{Plataforma} & \textbf{Destino} & \textbf{Alojamiento supuesto} & \textbf{Aporte del proponente} \\
\midrule\noalign{}
\endhead
\bottomrule\noalign{}
\multicolumn{4}{@{}p{\columnwidth}@{}}{\rule{0pt}{3.2ex}\textbf{Fuente:} Elaboración propia a partir del Subdocumento 4 (Tablas 4.7 y 4.22).} \\
\endlastfoot
ERP y emisión de documentos tributarios & Se conserva como único emisor tributario. & DC-01 & Adaptador y modalidad fiscal de contingencia. \\ \addlinespace
WMS del centro de distribución principal & Se conserva; su extensión a Concepción se evalúa. & CD-01 & Adaptador y degradación por interfaz. \\ \addlinespace
Marketplace de 2022 & Se conserva. & Externo & Adaptador del Servicio de marketplace (R:V-04). \\ \addlinespace
Comercio electrónico de 2019 & Se integra durante la evaluación de Etapa 1. & Externo & Adaptador y pruebas de conservación o sustitución. \\ \addlinespace
Fidelización de 2017 & Se integra durante la evaluación de Etapa 1. & Externo & Adaptador y pruebas de conservación o sustitución. \\ \addlinespace
Núcleo Retail de 2009 & Se reemplaza por olas conciliadas. & DC-01 & Servicios de Retail de la fila 2.9. \\ \addlinespace
Plataforma de crédito de 2011 & Se reemplaza por olas de cartera y funciones. & Ubicación por acreditar & Servicios del Emisor de la fila 2.9. \\ \addlinespace
Punto de venta de 2014 & Se sustituye tienda por tienda. & Terminales de SIT-01 & Filas 2.11 y 2.12. \\ \addlinespace
Motor de precios y promociones & No existe como sistema; es capacidad nueva. & No aplica & Servicio de oferta comercial (R:M-01). \\
\end{longtable}

La Tabla~\mbox{\ref{tab:t11-9}} delimita la oferta. El Cliente conserva las licencias y el hardware de las plataformas que se mantienen, y el proponente aporta solo los adaptadores y la integración. El alojamiento de cada plataforma se confirma con el inventario, y un cambio de ubicación modifica la ruta del adaptador respectivo sin alterar los servicios nuevos.

\subsection{T-11.9 Componentes de las innovaciones}\label{t-11-9}

Las cinco innovaciones del Subdocumento 13 agregan equipos y plataformas propios, que entran en producción en el mes 21, salvo el conteo predictivo, que opera desde el mes 16. La Tabla~\mbox{\ref{tab:t11-10}} los reúne para que la arquitectura, el plan y la oferta económica usen la misma lista.

{\footnotesize
\begin{longtable}[]{@{}
  >{\raggedright\arraybackslash}p{(\columnwidth - 10\tabcolsep) * \real{0.1500}}
  >{\raggedright\arraybackslash}p{(\columnwidth - 10\tabcolsep) * \real{0.1700}}
  >{\raggedright\arraybackslash}p{(\columnwidth - 10\tabcolsep) * \real{0.2000}}
  >{\raggedright\arraybackslash}p{(\columnwidth - 10\tabcolsep) * \real{0.1400}}
  >{\raggedright\arraybackslash}p{(\columnwidth - 10\tabcolsep) * \real{0.1000}}
  >{\raggedright\arraybackslash}p{(\columnwidth - 10\tabcolsep) * \real{0.2400}}@{}}
\caption{Componentes de las innovaciones.}\label{tab:t11-10}\\
\toprule\noalign{}
\textbf{Innovación} & \textbf{Componente} & \textbf{Producto de referencia} & \textbf{Ubicación} & \textbf{Cantidad} & \textbf{Justificación} \\
\midrule\noalign{}
\endfirsthead
\toprule\noalign{}
\textbf{Innovación} & \textbf{Componente} & \textbf{Producto de referencia} & \textbf{Ubicación} & \textbf{Cantidad} & \textbf{Justificación} \\
\midrule\noalign{}
\endhead
\bottomrule\noalign{}
\multicolumn{6}{@{}p{\columnwidth}@{}}{\rule{0pt}{3.2ex}\textbf{Fuente:} Elaboración propia a partir del Subdocumento 13.} \\
\endlastfoot
Lockers Ancoa 24/7 & Gabinete electrónico exterior con controlador. & Cleveron 355 u opción técnicamente equivalente. & Tiendas definidas en el Subdocumento 13 & Por fijar en el levantamiento de la Etapa 2 & Documentado para uso a la intemperie y con configuración modular para mercancía voluminosa. \\ \addlinespace
Conteo cíclico predictivo & Modelo supervisado de árboles con boosting. & scikit-learn, sobre la carga de la fila 1.23. & CLD-01 & Un servicio & No requiere hardware en tienda ni identificación por radiofrecuencia. \\ \addlinespace
Búsqueda vectorial multimodal & Modelo de representación de imagen y texto. & Jina CLIP v2, con licencia comercial. & CLD-01 & Un servicio & Admite consulta por imagen y por texto en español. \\ \addlinespace
Búsqueda vectorial multimodal & Índice vectorial. & Extensión pgvector sobre el PostgreSQL de Retail (fila 1.5). & CLD-01 & Una base & Reutiliza el motor transaccional y evita operar un motor adicional. \\ \addlinespace
Retail Media Network & Plataforma de publicidad en sitio. & Criteo u opción técnicamente equivalente, con adaptador propio. & Servicio externo, EXT-01 & Un servicio & Plataforma especializada frente a un desarrollo propio. \\ \addlinespace
Modo Tienda con fila virtual de probadores & Tótem de autoservicio. & Zebra KC50 de 22 pulgadas, con Ethernet. & Tiendas definidas en el Subdocumento 13 & Por fijar en el levantamiento de la Etapa 2 & Montaje accesible y administración empresarial; la fila corre en el nodo local durante un corte. \\ \addlinespace
Modo Tienda con fila virtual de probadores & Impresora de turnos. & Epson TM-T20III Ethernet. & Junto a cada tótem & Uno por tótem & Entrega el turno en papel cuando el flujo lo requiere. \\
\end{longtable}
}

La Tabla~\mbox{\ref{tab:t11-10}} agrega dos cargas a componentes ya listados. La fila virtual se ejecuta en el nodo local durante los cortes y entra en su dimensionamiento, y el índice vectorial usa el PostgreSQL de Retail, cuya admisión de la extensión se comprueba en la suscripción.

\subsection{T-11.10 Licenciamiento}\label{t-11-10}

Durante el contrato, el proponente contrata y opera las suscripciones de Azure y los servicios de la solución (Tabla 4.22 del Subdocumento 4) y las herramientas con que construye y entrega. Las licencias de los componentes que el Cliente usa al terminar el contrato se transfieren a su nombre. Todas rigen durante los 56 meses. La Tabla~\mbox{\ref{tab:t11-11}} declara la modalidad y la licencia de cada producto, incluido el software de código abierto y su compatibilidad con el uso previsto.

\begin{longtable}[]{@{}
  >{\raggedright\arraybackslash}p{(\columnwidth - 8\tabcolsep) * \real{0.2200}}
  >{\raggedright\arraybackslash}p{(\columnwidth - 8\tabcolsep) * \real{0.1700}}
  >{\raggedright\arraybackslash}p{(\columnwidth - 8\tabcolsep) * \real{0.1800}}
  >{\raggedright\arraybackslash}p{(\columnwidth - 8\tabcolsep) * \real{0.1100}}
  >{\raggedright\arraybackslash}p{(\columnwidth - 8\tabcolsep) * \real{0.3200}}@{}}
\caption{Licenciamiento por producto.}\label{tab:t11-11}\\
\toprule\noalign{}
\textbf{Producto} & \textbf{Modalidad} & \textbf{Licencia} & \textbf{Titular} & \textbf{Compatibilidad con el uso previsto} \\
\midrule\noalign{}
\endfirsthead
\toprule\noalign{}
\textbf{Producto} & \textbf{Modalidad} & \textbf{Licencia} & \textbf{Titular} & \textbf{Compatibilidad con el uso previsto} \\
\midrule\noalign{}
\endhead
\bottomrule\noalign{}
\multicolumn{5}{@{}p{\columnwidth}@{}}{\rule{0pt}{3.2ex}\textbf{Fuente:} Elaboración propia.} \\
\endlastfoot
Servicios de Azure de la Tabla~\mbox{\ref{tab:t11-2}} & Suscripción por consumo y por capacidad reservada. & Contrato de servicios en línea de Microsoft. & Proponente durante el contrato. Traspaso al Cliente según el Plan de Reversibilidad. & Uso comercial sin restricción. \\ \addlinespace
Microsoft Entra ID P2, Power BI, Sentinel y Defender & Suscripción por usuario o por volumen. & Contrato de servicios en línea de Microsoft. & Proponente durante el contrato. Traspaso al Cliente según el Plan de Reversibilidad. & Uso comercial sin restricción. \\ \addlinespace
Azure Pipelines, SonarQube, Jira, SharePoint y Slack & Suscripción por usuario. & Licencia comercial del fabricante. & Proponente. Herramientas de trabajo que no se transfieren, salvo el canal de entrega si el Cliente lo pide. & Uso comercial sin restricción. \\ \addlinespace
Java 25, distribución OpenJDK & Código abierto con soporte del distribuidor. & GPL versión 2 con excepción Classpath. & No requiere titular & La excepción permite ejecutar aplicaciones propias sin obligación de liberarlas. \\ \addlinespace
Spring Boot, TypeScript, Kubernetes, OpenTelemetry, Trivy, Syft, Cosign, OWASP ZAP y Playwright & Código abierto. & Apache 2.0. & No requiere titular & Permisiva; sin obligación de liberar el código propio. \\ \addlinespace
React y Pact & Código abierto. & MIT. & No requiere titular & Permisiva. \\ \addlinespace
PostgreSQL y pgvector & Código abierto, dentro del servicio administrado. & Licencia PostgreSQL. & No requiere titular & Permisiva. \\ \addlinespace
scikit-learn & Código abierto. & BSD de tres cláusulas. & No requiere titular & Permisiva. \\ \addlinespace
Terraform & Código disponible. & Business Source License 1.1. & No requiere titular & Admite el uso interno para administrar infraestructura propia; no se redistribuye. \\ \addlinespace
k6 & Código abierto. & AGPL versión 3. & No requiere titular & Se usa como herramienta de prueba, sin modificarla ni ofrecerla como servicio, por lo que no impone obligaciones sobre la solución. \\ \addlinespace
Jina CLIP v2 & Licencia comercial del modelo. & Licencia comercial. La versión pública es no comercial. & Cliente, desde el suministro & El suministro incluye la licencia comercial. \\ \addlinespace
Plataforma de publicidad y software de los casilleros & Suscripción o licencia del fabricante. & Licencia comercial del fabricante. & Cliente & Uso comercial sin restricción. \\ \addlinespace
Servicios de negocio, adaptadores, punto de venta y nodo local & Desarrollo propio. & Propiedad transferida al Cliente. & Cliente & No contiene componentes preexistentes del proponente. \\
\end{longtable}

La Tabla~\mbox{\ref{tab:t11-11}} muestra que ninguna licencia de código abierto impone la liberación del código de la solución. Las dos licencias que exigen atención son la del modelo multimodal, que se adquiere en su modalidad comercial, y la de la herramienta de carga, que se usa sin modificar. El registro de las licencias que se transfieren al Cliente es el paquete 1.4.7 de la EDT y se completa en el hito H3.

OpenSearch, OpenMetadata, MLflow, React Native y las herramientas de observabilidad se incorporan al registro de licencias por versión y modalidad de uso. Su compatibilidad no se deduce de la licencia de otros productos. El suministro debe comprobar las condiciones de uso, soporte y actualización de cada versión antes de su puesta en servicio.

\subsection{T-11.11 Ciclo de vida y soporte}\label{t-11-11}

Todo componente ofertado tiene soporte vigente del fabricante o de su comunidad. La Tabla~\mbox{\ref{tab:t11-12}} declara, para cada componente principal, la versión ofertada, el fin de soporte publicado y el plan de actualización durante el contrato.

\begin{longtable}[]{@{}
  >{\raggedright\arraybackslash}p{(\columnwidth - 6\tabcolsep) * \real{0.2200}}
  >{\raggedright\arraybackslash}p{(\columnwidth - 6\tabcolsep) * \real{0.1700}}
  >{\raggedright\arraybackslash}p{(\columnwidth - 6\tabcolsep) * \real{0.2900}}
  >{\raggedright\arraybackslash}p{(\columnwidth - 6\tabcolsep) * \real{0.3200}}@{}}
\caption{Versión, fin de soporte y plan de actualización.}\label{tab:t11-12}\\
\toprule\noalign{}
\textbf{Componente} & \textbf{Versión ofertada} & \textbf{Fin de soporte publicado} & \textbf{Plan de actualización en el contrato} \\
\midrule\noalign{}
\endfirsthead
\toprule\noalign{}
\textbf{Componente} & \textbf{Versión ofertada} & \textbf{Fin de soporte publicado} & \textbf{Plan de actualización en el contrato} \\
\midrule\noalign{}
\endhead
\bottomrule\noalign{}
\multicolumn{4}{@{}p{\columnwidth}@{}}{\rule{0pt}{3.2ex}\textbf{Fuente:} Elaboración propia a partir de las referencias citadas, consultadas en octubre de 2026.} \\
\endlastfoot
Java & 25 LTS. & Septiembre de 2030 en soporte principal y septiembre de 2033 en soporte extendido (Oracle, 2026). & Migración a la siguiente versión de soporte prolongado durante la operación, antes de septiembre de 2030. \\ \addlinespace
Spring Boot & Línea 4. Versión menor por fijar en el hito H3 (mes 6). & Cada versión menor recibe soporte comunitario por unos 13 meses; el de la versión 4.0 termina en diciembre de 2026 (Spring, 2026). & Actualización a cada versión menor dentro de su ventana de soporte, al menos una vez por año. \\ \addlinespace
Kubernetes en AKS & Versión admitida por el servicio, por fijar en el hito H3 (mes 6). & El servicio admite las tres versiones menores más recientes (Microsoft, 2026b). & Actualización programada al menos una vez por año, fuera de las ventanas de congelamiento del Caso. \\ \addlinespace
PostgreSQL & Versión mayor 17 o 18, por fijar en el hito H3 (mes 6). & El servicio sigue el calendario de la comunidad, que retira la versión 17 en noviembre de 2029 y la 18 en noviembre de 2030 (Microsoft, 2026c; PostgreSQL Global Development Group, 2026). & Una actualización de versión mayor durante la operación, ensayada en Preproducción. \\ \addlinespace
Azure Managed Redis & Versión administrada por el servicio. & Sin retiro anunciado; reemplaza a Azure Cache for Redis, que se retira el 30 de septiembre de 2028 (Microsoft, 2026a). & Actualización administrada por el proveedor. \\ \addlinespace
API Management Premium clásico, Service Bus Premium y demás servicios administrados & Versión administrada por el servicio. & Sin retiro anunciado a la fecha de consulta. & Revisión trimestral de los avisos de retiro del proveedor, con migración planificada con un año de anticipación (riesgo R-07 del Formulario T-16). \\ \addlinespace
TypeScript y React & Versiones estables vigentes. & Sin calendario fijo; soporte sobre la versión mayor vigente. & Actualización anual de componentes de base. \\ \addlinespace
Terraform y herramientas de la cadena de entrega, entre ellas Azure Pipelines & Versiones estables vigentes. & Según cada proyecto. & Actualización anual, con las versiones fijadas en el repositorio. \\ \addlinespace
Nodo local de tienda & Equipo nuevo, modelo por fijar en los meses 6 y 7. & Garantía y soporte del fabricante por los 56 meses, contratados por el Cliente. & Reposición adquirida por el Cliente si el modelo deja de tener soporte antes del mes 56, con especificación, configuración y certificación del reemplazo por el proponente. Borrado seguro certificado y disposición final con gestor autorizado al retiro. \\
\end{longtable}

La Tabla~\mbox{\ref{tab:t11-12}} identifica tres actualizaciones que ocurrirán con certeza dentro del contrato, que son las versiones menores de Spring Boot, las versiones de Kubernetes y una versión mayor de PostgreSQL, además de la migración de Java antes de 2030. Todas quedan a cargo del proponente como parte de la mantención, se ensayan en Preproducción y respetan las ventanas de congelamiento del Caso. El calendario de ciclo de vida se mantiene en la cuenta 1.15.4 de la EDT.

\subsection{T-11.12 Portabilidad}\label{t-11-12}

La estrategia de reversibilidad distingue los componentes que se trasladan a otro proveedor sin reescritura de los que son propios de Azure. La Tabla~\mbox{\ref{tab:t11-13}} clasifica cada grupo e indica en qué consiste su migración.

\begin{longtable}[]{@{}
  >{\raggedright\arraybackslash}p{(\columnwidth - 4\tabcolsep) * \real{0.3000}}
  >{\raggedright\arraybackslash}p{(\columnwidth - 4\tabcolsep) * \real{0.2000}}
  >{\raggedright\arraybackslash}p{(\columnwidth - 4\tabcolsep) * \real{0.5000}}@{}}
\caption{Portabilidad de los componentes.}\label{tab:t11-13}\\
\toprule\noalign{}
\textbf{Componente} & \textbf{Clasificación} & \textbf{Qué implica migrarlo} \\
\midrule\noalign{}
\endfirsthead
\toprule\noalign{}
\textbf{Componente} & \textbf{Clasificación} & \textbf{Qué implica migrarlo} \\
\midrule\noalign{}
\endhead
\bottomrule\noalign{}
\multicolumn{3}{@{}p{\columnwidth}@{}}{\rule{0pt}{3.2ex}\textbf{Fuente:} Elaboración propia a partir del Subdocumento 4 (Tablas 4.12 y 4.20).} \\
\endlastfoot
Servicios de negocio y adaptadores en contenedores & Portable & Desplegar las mismas imágenes en otro Kubernetes. \\ \addlinespace
PostgreSQL & Portable & Exportar e importar con herramientas estándar del motor. \\ \addlinespace
Clientes de Service Bus & Adaptación al transporte & Migrar contratos, sesiones, errores y deduplicación a las garantías del intermediario de destino. \\ \addlinespace
Instrumentación OpenTelemetry & Portable & Cambiar el destino de exportación. \\ \addlinespace
Datos del lago analítico & Portable & Copiar archivos en formatos abiertos. \\ \addlinespace
Software del punto de venta y del nodo local & Portable & No depende de la nube elegida. \\ \addlinespace
Definición de infraestructura en Terraform & Portable en la herramienta & Reescribir los módulos para el proveedor de destino. \\ \addlinespace
API Management & Propio del proveedor & Traducir políticas, productos y suscripciones a otra puerta de enlace. \\ \addlinespace
Service Bus y recuperación desde registros durables & Propio del proveedor & Migrar entidades y clientes, conservando registros de productores y consumidores para el reproceso. \\ \addlinespace
Front Door, Application Gateway, Azure Firewall y Bastion & Propio del proveedor & Reescribir reglas de borde y de red. \\ \addlinespace
Key Vault y Managed HSM & Propio del proveedor & Exportar las claves exportables y regenerar las que el hardware no libera. \\ \addlinespace
Entra ID & Propio del proveedor & Federar con otro proveedor de identidad mediante OIDC y SAML. \\ \addlinespace
Monitor, Log Analytics, Sentinel y Defender & Propio del proveedor & Exportar registros y reconstruir consultas y alertas. \\ \addlinespace
Power BI & Propio del proveedor & Reconstruir los informes sobre los mismos datos. \\
\end{longtable}

La Tabla~\mbox{\ref{tab:t11-13}} muestra que el código de negocio, los datos y el software de tienda son portables, y que la dependencia del proveedor se concentra en el borde, la identidad, las claves y la observabilidad. El esfuerzo de migración de cada componente propio del proveedor se estima en el Plan de Reversibilidad, y el riesgo de bloqueo se trata como R-08 en el Formulario T-16.

\subsection{T-11.13 Comprobaciones previas a la construcción}\label{t-11-13}

El Subdocumento 4 exige una comprobación de aceptación para cada tecnología antes de construir. La Tabla~\mbox{\ref{tab:t11-14}} reúne las comprobaciones que fijan los niveles, las cantidades y los modelos de este formulario, con el paquete de la EDT que las ejecuta y su hito.

\begin{longtable}[]{@{}
  >{\raggedright\arraybackslash}p{(\columnwidth - 6\tabcolsep) * \real{0.3000}}
  >{\raggedright\arraybackslash}p{(\columnwidth - 6\tabcolsep) * \real{0.3000}}
  >{\raggedright\arraybackslash}p{(\columnwidth - 6\tabcolsep) * \real{0.2400}}
  >{\raggedright\arraybackslash}p{(\columnwidth - 6\tabcolsep) * \real{0.1600}}@{}}
\caption{Comprobaciones que fijan cantidades, niveles y modelos.}\label{tab:t11-14}\\
\toprule\noalign{}
\textbf{Comprobación} & \textbf{Qué fija en este formulario} & \textbf{Paquete o instancia} & \textbf{Hito} \\
\midrule\noalign{}
\endfirsthead
\toprule\noalign{}
\textbf{Comprobación} & \textbf{Qué fija en este formulario} & \textbf{Paquete o instancia} & \textbf{Hito} \\
\midrule\noalign{}
\endhead
\bottomrule\noalign{}
\multicolumn{4}{@{}p{\columnwidth}@{}}{\rule{0pt}{3.2ex}\textbf{Fuente:} Elaboración propia a partir del Subdocumento 4 (Tablas 4.12, 4.15, 4.20 y 4.22) y del Formulario T-14.} \\
\endlastfoot
Creación efectiva de cada recurso en Chile Central y en Brazil South. & Niveles de las filas 1.2 a 1.31, incluidas la admisión de pgvector y la disponibilidad de Managed HSM (fila 1.24) y de Private Link para Front Door (fila 1.14). & Verificación del socio certificado de Microsoft en la suscripción del proyecto. & Antes de H3, mes 6 \\ \addlinespace
Memoria de capacidad y diseño detallado. & Nodos de AKS, capacidad de Service Bus y registros durables de recuperación, clase de cómputo de PostgreSQL, tamaño de la caché, grupos de nodos predictivos, cantidades de las filas 1.25 a 1.30, estaciones de la fila 5.13 y ancho de banda por enlace. & Dimensionamiento de la sección 4.2. & H2, mes 4 \\ \addlinespace
Inventario por tienda y levantamiento de enlaces. & Cantidades finales de las filas 5.1 a 5.6, 5.11 y 5.14, modelos de referencia de la red de sitio, y enlaces 6.2 y 6.3. & Paquete 1.4.8 de la EDT. & Meses 3 a 6 \\ \addlinespace
Inventario de las catorce interfaces y de la ubicación de cada plataforma. & Ubicación de los adaptadores de la fila 2.10 y necesidad del enlace 6.7. & Levantamiento de integración. & Antes de H3, mes 6 \\ \addlinespace
Prototipo del punto de venta. & Empaquetado de la fila 2.11 y compatibilidad de las filas 5.7 y 5.8. & Prototipo de caja. & Antes del piloto, mes 6 \\ \addlinespace
Ensayo de continuidad de 24 horas en una tienda piloto. & Modelo de referencia, capacidad, consumo, margen de crecimiento, motor, empaquetado y sistema operativo del nodo local, y potencia de la fila 5.4. & Piloto de corte y reconexión. & Meses 6 y 7 \\ \addlinespace
Prueba de permisos y linaje. & Modelos PostgreSQL analíticos, modalidad Power BI por informe y permisos de búsqueda Retail y catálogo. & Prueba de analítica por negocio. & Antes de producción de la etapa correspondiente. \\ \addlinespace
Aprobación de residencia por ámbito. & Destino aprobado por ámbito y alternativa nacional si Brasil se rechaza. & Cliente y proponente, con revisión jurídica y ensayo técnico. & Aprobación antes del cierre de diseño y ensayo antes de producción. \\ \addlinespace
Informe interno de 2024 del centro de datos. & Alcance del plan de cierre de brecha de DC-01 y cantidades de la Tabla~\mbox{\ref{tab:t11-dc01}}. & Paquetes 1.4.9 a 1.4.13 de la EDT. & H3, mes 6 \\
\end{longtable}

Las comprobaciones distinguen la selección de diseño de los resultados que deben obtenerse en ensayo. El inventario se actualiza bajo control de cambios y se concilia con las Tablas 4.20 y 4.22 del SD-04. La viabilidad regional de API Management, el reproceso durable y la capacidad de observabilidad de X-01 deben verificarse antes de su aceptación.

\subsection{Referencias}\label{t11-referencias}

Las siguientes fuentes respaldan los datos del Caso y las fechas de soporte citados en este formulario.

\noindent\hangindent=1.5em\hangafter=1 Microsoft. (2026a). \emph{What's new in Azure Cache for Redis}. \url{https://learn.microsoft.com/azure/azure-cache-for-redis/cache-whats-new}

\noindent\hangindent=1.5em\hangafter=1 Microsoft. (2026b). \emph{Supported Kubernetes versions in Azure Kubernetes Service (AKS)}. \url{https://learn.microsoft.com/azure/aks/supported-kubernetes-versions}

\noindent\hangindent=1.5em\hangafter=1 Microsoft. (2026c). \emph{Azure Database for PostgreSQL versioning policy}. \url{https://learn.microsoft.com/azure/postgresql/configure-maintain/concepts-version-policy}

\noindent\hangindent=1.5em\hangafter=1 Multitiendas Ancoa S.A. (2026b). \emph{Bases técnicas. Caso 09: Cadena multitienda} (Licitación N.° TFEP-01/2026, versión 1.0).

\noindent\hangindent=1.5em\hangafter=1 Oracle. (2026). \emph{Java SE Support Roadmap}. \url{https://www.oracle.com/java/technologies/java-se-support-roadmap.html}

\noindent\hangindent=1.5em\hangafter=1 PostgreSQL Global Development Group. (2026). \emph{Versioning policy}. \url{https://www.postgresql.org/support/versioning/}

\noindent\hangindent=1.5em\hangafter=1 Spring. (2026). \emph{Spring Boot support}. \url{https://spring.io/projects/spring-boot#support}

\ossFinalPage

\end{document}
